\subsection*{Learning where to find an answer}
A wrong pointer can produce legible words that do not answer the question. We tested two sources of supervision for the next written word: answers generated by a local 4B text model, and supporting passages identified by human readers. The text teacher answered from a passage, a current question and the previous human exchange. Among 83 turns where all four references required absent multiword phrasing, it matched a human answer on ten and obtained 21.1\% official word F1. It returned yes, no or unknown on 58. Placing the current question last raised the exact count to 13 and F1 to 27.6\%. On nine separate difficult turns from typographically normalized stories, F1 was 32.0\%. Both tests fell short of the prespecified 35\% required for teacher supervision. The student fits consequently used the human responses and evidence intervals alone (Fig.~\ref{fig:dialogue-evidence}).

We added a loss that rewards the total pointer probability of all correct word images inside the annotated interval. On 758 turns from 47 previously excluded development stories of 321--448 words, the visual writer reached 41.84\% official word F1 after this continuation, compared with 41.08\% after the same additional training without the new loss. The paired difference was 0.76 percentage points (95\% whole-story interval 0.003 to 1.60). Corresponding byte scores were 42.57\% and 42.47\%, a difference of 0.10 points ($-0.79$ to $+0.98$). On a question asking who first took tea, the visual answer changed from \emph{Black} to \emph{Shen Nong}. A different question about Cockney speech changed from a correct name to \emph{Madam}. Neither visual arm matched a reference on the 22 turns where all four human answers required absent wording. The visual gain is small and sensitive to which stories are drawn.

Across four paired visual fits from the same initial checkpoint, the mean effect was +0.39 percentage points on the 47 stories (95\% two-level interval $-0.23$ to $+1.06$) and +0.27 on eight longer stories ($-0.69$ to $+1.47$). On the longer stories two seeds had positive and two had negative differences. Across all eight visual fits, none matched any reference on the 27 turns requiring new multiword phrasing. The original improvement did not meet the criterion for reliable transfer across training batches and stories.

Giving the parent model the annotated interval during a separate diagnostic raised the mean probability of the required copyable word from 0.398 to 0.502. Upweighting the interval corrected 135 top-location decisions and spoiled ten among 1,834 positions under correct previous words. The autonomous writer can instead estimate membership of an interval from its own endpoint scores. This replacement changed visual F1 from 41.84\% to 41.97\% on the 758 turns, and from 35.67\% to 36.08\% on another 154 turns from eight longer stories. Both paired intervals included zero. Evidence supplied by a human can identify the requested content more reliably than the measured improvement achieved by this estimated-interval decoder.

All 27 difficult fresh turns admitted at least one normalized human answer within the writer's action space. Eighteen could be answered entirely by rearranging visible word images. The other nine required spelling only one word alongside any copied words. Of the 108 human references, 59 admitted an entirely copied answer and 107 could be written by some legal action sequence. We established reachability by searching over the four physical edge edits per observed word, the fitted letter palette and the model's termination limits, and verified each recovered sequence against the corresponding reference. The model's zero exact answers therefore localize a failure in choosing and ordering content even when the required glyphs are available.
